% UL_Report.tex — CS7641 UL Report (Summer 2026)
% Timothy Leung (tleung37) | gtID 904150400 | Git seed 7641
% Style, palette, macros identical to SL_Report.tex and OL_Report.tex

\documentclass[conference]{IEEEtran}
\IEEEoverridecommandlockouts

% ── packages ──────────────────────────────────────────────────────────────────
\usepackage{cite}
\usepackage{amsmath,amssymb,amsfonts}
\usepackage{graphicx}
\usepackage{textcomp}
\usepackage[table]{xcolor}
\usepackage{booktabs}
\usepackage{multirow}
\usepackage{tabularx}
\usepackage{array}
\usepackage{caption}
\usepackage{subcaption}
\usepackage[hidelinks]{hyperref}
\usepackage{enumitem}
\usepackage{fancyhdr}

% ── denim palette ─────────────────────────────────────────────────────────────
\definecolor{denim}{HTML}{2C6FBB}
\definecolor{denimdark}{HTML}{1B3A5C}
\definecolor{denimmid}{HTML}{3A6B8C}
\definecolor{denimbright}{HTML}{5B9BD5}
\definecolor{denimrow}{HTML}{E6EEF5}
\definecolor{amber}{HTML}{D4A017}
\definecolor{plum}{HTML}{7B2D8E}
\definecolor{sage}{HTML}{6BB38A}
\definecolor{sienna}{HTML}{B8552E}
\definecolor{teal}{HTML}{4C9E82}
\definecolor{parchment}{HTML}{F4EFE5}
\definecolor{denim_comp}{HTML}{BB782C}
\definecolor{plum_comp}{HTML}{408E2D}

\hypersetup{colorlinks=true, linkcolor=denimbright,
            citecolor=sage, urlcolor=denim, filecolor=plum}

% ── running header ────────────────────────────────────────────────────────────
\pagestyle{fancy}
\fancyhf{}
\fancyhead[L]{\small\color{denimmid}CS7641 UL Report --- Timothy Leung}
\fancyhead[R]{\small\color{denimmid}gtID 904150400 \textbar\ Git seed 7641}
\fancyfoot[C]{\small\color{denimmid}\thepage}
\renewcommand{\headrulewidth}{0.3pt}

% ── macros ────────────────────────────────────────────────────────────────────

\newcommand{\figref}[1]{\hyperref[#1]{\textcolor{denim_comp}{Fig.~\ref*{#1}}}}
\newcommand{\tabref}[1]{\hyperref[#1]{\textcolor{sage}{Table~\ref*{#1}}}}
\newcommand{\secref}[1]{\hyperref[#1]{\textcolor{plum}{\S\ref*{#1}}}}
\newcommand{\backto}[1]{\hfill\hyperref[#1]{\textcolor{denimmid}{\footnotesize $\hookleftarrow$ back}}}
\newcommand{\partref}[2]{\hyperref[#2]{\textcolor{teal}{Part~#1}}}

% ── listings style ────────────────────────────────────────────────────────────
\usepackage{listings}
\lstset{basicstyle=\ttfamily\footnotesize,
        breaklines=true, frame=none,
        keywordstyle=\color{denim},
        commentstyle=\color{sage}}

% ══════════════════════════════════════════════════════════════════════════════
\begin{document}
\graphicspath{{../results/figures/}{./results/figures/}}

\title{\color{denimdark}Unsupervised Structure in Covertype:\\
       Clustering, Dimensionality Reduction, and Downstream Learning}

\author{
  \IEEEauthorblockN{Timothy Leung}
  \IEEEauthorblockA{Georgia Institute of Technology, OMSCS\\
  CS 7641 --- Machine Learning, Summer 2026\\
  gtID 904150400 \quad \texttt{tleung37@gatech.edu}}
}

\maketitle

% ── Abstract ──────────────────────────────────────────────────────────────────
\begin{abstract}
We investigate the geometric structure of the Forest Cover Type (Covertype)
dataset through unsupervised learning. K-Means and Gaussian Mixture Models
(GMM/EM) applied to the original 54-dimensional feature space yield low
cluster--label alignment (ARI $\leq 0.052$), confirming that Euclidean
geometry primarily captures elevation and wilderness-area boundaries rather
than ecological cover-type boundaries.
Principal Component Analysis (PCA) retains 90\,\% of variance in only 10
components, ICA extracts 20 high-kurtosis independent components
(mean $|\kappa|{=}6.05$), and Randomized Projections (RP) preserve mean pairwise distances
faithfully across five seeds (seed-to-seed variation in mean ratio: $0.0045$),
but individual pair distortions reach ${\pm}12\%$ of original distance.
PCA-reduced features yield the highest downstream MLP Macro-F1 of
$\mathbf{0.6951}$, outperforming both the original-feature baseline
(0.6803) and ICA/RP representations.
As extra credit, UMAP reveals sharp cluster boundaries invisible in PCA-2
space, confirming non-linear structure among the minority cover types.
\end{abstract}

% ── Section I: Introduction ───────────────────────────────────────────────────
\section{Introduction}\label{sec:intro}

This report is the third in the Covertype sequence following the Supervised
Learning (SL) and Optimization \& Uncertainty (OL) Reports.
The SL Report established that a PyTorch MLP [256, 128] with Adam and L2
regularization achieves macro-F1~$\approx 0.69$ on the held-out test set.
The OL Report identified Adam without bias correction (Adam-NB) as the best
optimizer and L2 ($\lambda{=}10^{-4}$) as the best single regularizer.
This report asks the complementary question: \emph{what structure exists in
the feature space before labels are used, and do lower-dimensional
representations help or hurt the best supervised workflow?}

\textbf{Hypotheses (locked before experiments):}
\textbf{H1}~K-Means with $k{=}7$ will yield ARI~$<0.15$, as Euclidean
geometry captures elevation/wilderness structure rather than ecological labels.
\textbf{H2}~PCA will preserve downstream MLP macro-F1 within 0.02 of the
original-feature baseline at $\leq 30$ components, outperforming ICA and RP
because Covertype's correlated cartographic features align with principal
variance axes.
\textbf{H3}~PCA will modestly improve K-Means silhouette over original space
by suppressing sparse binary indicator noise; RP will show the highest
seed-to-seed ARI variability.
A failure of any clause is informative; the verdict per clause is recorded
in \secref{sec:conclusion}.

\secref{sec:data} fixes the Covertype continuity contract and metric
justification; \secref{sec:p1} through \secref{sec:p4} report Parts~1--4;
\secref{sec:umap} covers the extra-credit UMAP visualization;
\secref{sec:disc} synthesises across all four parts;
\secref{sec:conclusion} closes with three direct answers to the three
opening questions.

% ── Section II: Data and Preprocessing ───────────────────────────────────────
\section{Data and Preprocessing}\label{sec:data}

The working dataset is the same stratified 20,000-instance Covertype sample
from the SL Report (seed~7641, \texttt{StratifiedShuffleSplit}).
The 54-feature matrix comprises 10 continuous cartographic variables (columns
0--9: elevation, aspect, slope, distance-to-hydrology, hillshade, and
distance-to-fire-points) and 44 binary wilderness-area/soil-type indicators
(columns 10--53).

\textbf{Preprocessing decision.}
Continuous columns are standardised (\texttt{StandardScaler} fit on training
split only for downstream NN; fit on full 20k for exploratory clustering).
Binary columns are left as 0/1 without rescaling: scaling them would inflate
their relative variance contribution to PCA and inflate distances for K-Means
disproportionately~\cite{jolliffe2002pca}.

\textbf{What the deliberate stratification implies for interpretation.}
The 20k stratified sample preserves CT4 Cottonwood/Willow at $0.47\%$
(${\approx}94$ rows total, ${\approx}19$ test instances).
This is a deliberate choice with a cost: 19 test instances means each per-class
metric for CT4 carries a ${\pm}0.10$ standard error on recall and
${\pm}0.10$ on precision.
That cost is worth paying because the alternative --- oversampling or
dropping CT4 --- would destroy the ecological signal that motivates the
unsupervised analysis.
The question this report is really asking is: \emph{does the feature space
contain exploitable structure that corresponds to the seven ecological niches,
or is the data fundamentally a cartographic problem where the biology is
beyond the sensor's resolution?}
Answering that question requires CT4 to be present in the analysis, even if
every per-class number for CT4 must be read as approximate.
Macro-F1 is used as the downstream NN metric because its equal per-class
weighting means CT4's ${\approx}19$ test instances carry the same
contribution to the reported score as CT1's ${\approx}1{,}459$ ---
making it sensitive to whether dimensionality reduction helps or hurts
the ecologically rarest and geometrically most isolated class.
This choice yields 7 class proportions matching the SL Report exactly
(see \figref{fig:class_dist}).

The 60/20/20 train/val/test split is identical to SL and OL.
For exploratory clustering and DR diagnostics (Parts~1--2), we use the full
20k scaled sample.
For downstream NN evaluation (\partref{4}{sec:p4}), all DR transformations are fit on the
training split only and applied to validation/test~\cite{sklearn2024}.

\begin{figure}[!t]
  \centering
  \includegraphics[width=\linewidth]{./results/figures/fig_class_distribution.pdf}
  \caption{Class distribution confirms the same imbalanced Covertype sample as
  SL/OL Reports. CT1 (42\,\%) and CT2 (49\,\%) dominate; CT4--CT7 are minority
  classes that challenge both clustering and supervised learning.}
  \label{fig:class_dist}\backto{sec:data}
\end{figure}

% ── Section III: Part 1 — Clustering on Original Features ────────────────────
\section{Part 1: Clustering on Original Features}\label{sec:p1}

\subsection{Parameter justification}\label{sec:p1:params}

\textbf{K-Means.} We sweep $k \in \{2,3,4,5,6,7,8,10,12\}$ with
\texttt{n\_init=20} and \texttt{max\_iter=500} (see \figref{fig:clust_sweep}).
The inertia elbow is diffuse; no sharp breakpoint at $k{=}7$.
Silhouette peaks at $k{=}2$ (0.190) and declines monotonically, suggesting the
data lacks compact spherical clusters at any granularity.
We fix $k{=}7$ as the \emph{reference} matching the known label count, enabling
fair label-alignment comparison across Parts~1 and~3.

\textbf{GMM/EM.}
We sweep $n \in \{2,3,\ldots,12\}$ with full covariance and \texttt{n\_init=5}.
BIC continues to decrease through $n{=}12$, consistent with the overlapping,
non-spherical class geometry suggested by the PCA analysis.
We fix $n{=}7$ for direct comparison with K-Means and Cover~Type.

\begin{figure}[!t]
  \centering
  \includegraphics[width=\linewidth]{./results/figures/fig_clustering_sweep.pdf}
  \caption{\emph{Left}: inertia and BIC sweeps --- no sharp elbow at $k{=}7$.
  \emph{Center}: silhouette is highest at $k{=}2$ for K-Means; GMM silhouette is
  consistently lower, indicating its elliptical components do not match
  distance-based silhouette geometry.
  \emph{Right}: ARI vs.\ Cover~Type labels is low across all $k$, peaking at
  $k{=}3$ for GMM (ARI~$=$~0.119), consistent with the two dominant classes
  (CT1, CT2) and an elevation-driven third grouping.
  Vertical \textcolor{sienna}{\emph{dotted}} line = $k{=}7$ (reference).}
  \label{fig:clust_sweep}\backto{sec:p1:params}
\end{figure}

\subsection{Results and label alignment}\label{sec:p1:results}

\tabref{tab:clust_orig} summarises the final metrics.
K-Means ARI of 0.023 and NMI of 0.067 confirm \textbf{H1}: Euclidean clusters
do not align with ecological labels.
The contingency analysis (\figref{fig:contingency}) shows that the two largest
K-Means clusters absorb nearly all CT1 and CT2 samples, while minority classes
CT3--CT7 are fragmented across clusters.
GMM achieves slightly higher NMI (0.131) because full-covariance components can
better capture the elongated CT1/CT2 manifolds, but ARI remains low (0.052).
Both algorithms primarily reflect the elevation gradient and
wilderness-area boundaries rather than the ecological taxonomy.
See \secref{sec:disc:hyp} for hypothesis verdict and \secref{sec:disc:cm}
for how this elevation-band structure reappears in the supervised confusion
matrices.

\begin{table}[!t]
\caption{\partref{1}{sec:p1} clustering metrics on original features ($n{=}20\,000$).
  \textbf{Best} of the two algorithms bolded per metric.}
\label{tab:clust_orig}\backto{sec:p1:results}
\centering\small
\rowcolors{2}{denimrow}{white}
\begin{tabular}{lrrrrr}
\toprule
\rowcolor{denimdark}
\textcolor{white}{\textbf{Algo}} & \textcolor{white}{\textbf{Sil}} & \textcolor{white}{\textbf{DB}} & \textcolor{white}{\textbf{CH}} &
\textcolor{white}{\textbf{ARI}} & \textcolor{white}{\textbf{NMI}} \\
\midrule
K-Means ($k{=}7$) & \textbf{0.1306} & \textbf{1.792} & \textbf{2683} & 0.023 & 0.067 \\
GMM (full, $n{=}7$) & $-$0.007 & 4.087 & 513 & \textbf{0.052} & \textbf{0.131} \\
\bottomrule
\end{tabular}
\end{table}

\begin{figure}[!t]
  \centering
  \includegraphics[width=\linewidth]{./results/figures/fig_contingency.pdf}
  \caption{Cluster $\times$ Cover~Type contingency (row-normalised) for K-Means
  (\emph{left}) and GMM (\emph{right}) at $n{=}7$.
  K-Means clusters are dominated by CT1/CT2 with minority types spread thinly;
  GMM distributes more across clusters but still fails to isolate CT4 (rarest class).}
  \label{fig:contingency}\backto{sec:p1:results}
\end{figure}

% ── Section IV: Part 2 — Dimensionality Reduction ────────────────────────────
\section{Part 2: Dimensionality Reduction}\label{sec:p2}

\subsection{PCA}\label{sec:p2:pca}

Full PCA (\figref{fig:pca_var}) shows that only \textbf{10 components} are
needed for 90\,\% cumulative variance and 14 for 95\,\%, verified identically
by R's \texttt{prcomp}.
The top two components account for 22.6\,\% and 19.7\,\% respectively,
consistent with the elevation ($\text{PC}_1$) and hillshade/aspect
($\text{PC}_2$) dominance.
The effective rank of the full 54-feature standardised matrix is $\approx 7.4$,
indicating strong covariance structure: 54 nominally independent features
behave as ${\approx}7$ independent signals.
(Restricting PCA to the 10 continuous features alone gives effective rank
$\approx 5.93$, as computed in the SL Report's factor analysis.
The difference reflects the binary block's contribution: 44 Bernoulli
dimensions add a small number of additional near-orthogonal axes that expand
the effective rank slightly beyond the continuous-only estimate.)
We fix \textbf{$n_{\text{PCA}}{=}30$} to capture 99.5\,\% of variance with a
stable, over-determined embedding for downstream NN; this exceeds the 90\,\%
target but avoids information loss at minimal compute cost~\cite{jolliffe2002pca}.

\begin{figure}[!t]
  \centering
  \includegraphics[width=\linewidth]{./results/figures/fig_pca_variance.pdf}
  \caption{\emph{Left}: individual EVR for top-20 PCs --- components 1--2 dominate,
  components beyond 14 contribute $<1\%$ each.
  \emph{Right}: cumulative EVR with 90\,\% (\emph{amber dashed}) and 95\,\% (\emph{sienna dashed})
  thresholds. $n{=}30$ (\emph{plum dotted}) captures 99.5\,\%, providing a generous
  but stable embedding.}
  \label{fig:pca_var}\backto{sec:p2:pca}
\end{figure}

\subsection{ICA}\label{sec:p2:ica}

FastICA with $n{=}20$ components and whitening yields a mean absolute kurtosis
of \textbf{6.05} across components (max~14.22), confirming non-Gaussian
independent source signals (\figref{fig:ica_kurt}).
The highest-kurtosis components likely capture soil-type indicator combinations
that are sparse and leptokurtic after whitening.
Seed sensitivity across 5 seeds shows mean absolute kurtosis is nearly
constant per seed ($6.01$--$6.08$, \figref{fig:ica_kurt} right) --- the
\emph{total non-Gaussianity} is stable across restarts --- but the
per-component kurtosis standard deviation across seeds is 4.58, meaning
individual components exchange orderings and sign-flip between runs.
This is the ICA instability that matters for reproducibility: the aggregate
signal is preserved, but component assignments are not,
a known limitation of FastICA on mixed continuous/binary
inputs~\cite{hyvarinen2000ica}.
We select \textbf{$n_{\text{ICA}}{=}20$} as the largest count that converges
reliably and provides a non-redundant representation.
A subtle point about what ICA's truncation discards: ICA orders components
by kurtosis (non-Gaussianity), so the 34 dropped components are those with
the \emph{smallest} absolute kurtosis.
On Covertype, those are the components encoding binary indicator combinations,
because Bernoulli marginals are already maximally non-Gaussian and ICA
cannot rank-order them further --- they saturate the kurtosis measure.
The components ICA selects preferentially are those driven by the continuous
cartographic features, which happen to carry the ecologically informative
signal.
This is why ICA's downstream MLP performance ($-0.002$) is near-parity
with the original: ICA's truncation roughly preserves the signal,
but its component instability across seeds prevents it from consistently
doing better than PCA's variance-ordered truncation.

\begin{figure}[!t]
  \centering
  \includegraphics[width=\linewidth]{./results/figures/fig_ica_kurtosis.pdf}
  \caption{\emph{Left}: per-component $|\text{kurtosis}|$ (excess) for 20 ICA components.
  All exceed the \textcolor{amber}{\emph{amber dashed}} $|\kappa|{=}3$ threshold, confirming non-Gaussian sources.
  \emph{Right}: mean $|\text{kurtosis}|$ per seed is nearly constant ($6.01$--$6.08$),
  confirming that total non-Gaussianity is stable across restarts.
  Seeds are plotted as unordered labels (no temporal structure implied).
  The ICA instability is in component \emph{ordering} and sign-flips across
  seeds, not in aggregate kurtosis
  --- see \secref{sec:p2:ica} for discussion.}
  \label{fig:ica_kurt}\backto{sec:p2:ica}
\end{figure}

\subsection{Randomized Projections}\label{sec:p2:rp}

Gaussian RP with $n{=}30$ preserves pairwise distances with mean ratio
$0.998 \pm 0.121$ across five seeds (\figref{fig:rp_dist}).
The mean ratio is near 1.0 (consistent with the Johnson--Lindenstrauss
lemma~\cite{johnson1984jl}), but the $\pm 12\,\%$ spread indicates
individual pair distances can be noticeably distorted.
Seed-to-seed variation in mean ratio is only 0.0045, suggesting that the
\emph{average} preservation is stable; individual pair distortions are the
main concern for distance-sensitive downstream methods (K-Means, GMM).

\begin{figure}[!t]
  \centering
  \includegraphics[width=\linewidth]{./results/figures/fig_rp_distortion.pdf}
  \caption{RP per-pair distance ratio distribution (violin + CDF, 500-point subsample,
  $n{=}30$ projected dimensions).
  \emph{Left}: all five seeds produce near-identical distributions centred at 1.0
  (JL lemma satisfied on average), but fat tails stretching from $0.5$ to $1.6$
  reveal that individual pairs are noticeably distorted.
  The seed-to-seed similarity of the violins confirms that mean-ratio variance
  (0.0045) is negligible --- falsifying \textbf{H3}'s seed-variability claim.
  \emph{Right}: CDF across seeds shows ${\approx}20\%$ of pairs fall outside
  the $[0.8,\,1.2]$ corridor (\textcolor{sienna}{\emph{dotted}} bounds) --- enough distortion to shift
  K-Means centroids and widen GMM covariance estimates in projected space,
  explaining RP's $-0.013$ Macro-F1 deficit vs.\ original features
  (\tabref{tab:nn_dr}).}
  \label{fig:rp_dist}\backto{sec:p2:rp}
\end{figure}

% ── Section V: Part 3 — Clustering after DR ───────────────────────────────────
\section{Part 3: Clustering after DR}\label{sec:p3}

\tabref{tab:clust_dr} reports absolute metrics; \figref{fig:dr_clust}
shows the $\Delta$ vs.\ original space, making the direction and magnitude
of each DR method's effect legible at a glance.

\textbf{K-Means} (\figref{fig:dr_clust}, left): all three DR methods improve
silhouette and NMI modestly (RP largest: $\Delta$sil~$=$~$+0.017$,
$\Delta$NMI~$=$~$+0.014$), partially supporting \textbf{H3}.
The outlier is ICA on ARI ($-0.028$): FastICA de-correlates the elevation
signal that K-Means centroids depend on, breaking dominant-variance geometry
without offering a distance-preserving replacement.

\textbf{GMM} (\figref{fig:dr_clust}, right): DR responses diverge sharply
between methods, revealing a mechanism-level tension.
ICA delivers the largest silhouette gain ($+0.137$) because whitening
sphericalises the input, giving GMM's full-covariance estimator a
numerically clean baseline --- the resulting components are
sphere-like, and GMM full-covariance can exploit that structure.
PCA dominates on ARI ($+0.065$) and NMI ($+0.044$): compressing 54 to 30
dimensions means GMM must estimate a $30{\times}30$ covariance matrix
instead of a $54{\times}54$ one --- with ${\approx}2{,}860$ points per
component at $n{=}7$, the 30-dim estimate is far more reliable than the
54-dim one, which is numerically rank-deficient and prone to covariance
collapse.
K-Means does not share this benefit because it never estimates a covariance
matrix --- its Euclidean centroid geometry is equally valid in 54 or 30
dimensions --- which is why K-Means silhouette improves only modestly
under PCA ($+0.002$) while GMM ARI improves by $+0.065$.
RP improves NMI ($+0.078$) but hurts ARI ($-0.025$) --- the same
distance distortion that benefits K-Means silhouette destroys the
dominant-variance geometry GMM needs for label-aligned components.
This algorithm $\times$ DR asymmetry is the core empirical finding of
\partref{3}{sec:p3}: no single DR method is universally best; the choice depends on
whether the downstream algorithm is distance-sensitive (K-Means) or
covariance-sensitive (GMM).

\begin{table}[!t]
\caption{Clustering metrics before and after DR.
  $\uparrow$~best silhouette; $\star$~best ARI per algorithm.}
\label{tab:clust_dr}\backto{sec:p3}
\centering\small
\rowcolors{2}{denimrow}{white}
\begin{tabular}{llrrrr}
\toprule
\rowcolor{denimdark}
\textcolor{white}{\textbf{Space}} & \textcolor{white}{\textbf{Algo}} & \textcolor{white}{\textbf{Sil}} & \textcolor{white}{\textbf{DB}} &
\textcolor{white}{\textbf{ARI}} & \textcolor{white}{\textbf{NMI}} \\
\midrule
Original & K-Means & 0.1306 & 1.792 & 0.023 & 0.067 \\
PCA-30   & K-Means & 0.1326 & 1.782 & 0.022 & 0.066 \\
ICA-20   & K-Means & 0.1416 & 1.701 & $-$0.005 & 0.105 \\
RP-30    & K-Means$\uparrow$ & \textbf{0.1478} & 1.650 & 0.027 & 0.080 \\
\midrule
Original & GMM     & $-$0.007 & 4.087 & 0.052 & 0.131 \\
PCA-30   & GMM$\star$ & $-$0.001 & 3.501 & \textbf{0.117} & 0.186 \\
ICA-20   & GMM     & 0.1304 & 1.780 & 0.027 & 0.143 \\
RP-30    & GMM     & 0.013 & 3.402 & 0.125 & 0.170 \\
\bottomrule
\end{tabular}
\end{table}

\begin{figure}[!t]
  \centering
  \includegraphics[width=\linewidth]{./results/figures/fig_dr_clustering_comparison.pdf}
  \caption{$\Delta$ clustering metrics after DR vs.\ original-space baseline
  (positive $=$ improvement). \emph{Left~(K-Means):} all three DR methods
  improve silhouette and NMI slightly; ICA hurts ARI ($-$0.028) because
  independent components disrupt the elevation-driven Euclidean structure
  K-Means relies on.
  \emph{Right~(GMM):} ICA delivers the largest silhouette gain ($+$0.137)
  because whitening regularises full-covariance estimation; PCA dominates
  on ARI ($+$0.065) and NMI ($+$0.044); RP improves NMI ($+$0.078) but
  hurts ARI ($-$0.025) by destroying dominant-variance geometry.
  The asymmetry between K-Means and GMM responses directly tests H3.}
  \label{fig:dr_clust}\backto{sec:p3}
\end{figure}

% ── Section VI: Part 4 — NN after DR ─────────────────────────────────────────
\section{Part 4: Neural Network after DR}\label{sec:p4}

We rerun the best-justified MLP from the SL and OL Reports
(architecture~$[256,128]$, Adam-NB optimizer, L2~$\lambda{=}10^{-4}$,
early stopping patience~10, 40 epochs, batch~256, seed~7641 ---
recipe derivation in \secref{sec:data}).
DR transformations are fit on the training split only and applied to
val/test.
\tabref{tab:nn_dr} reports overall test metrics.
\figref{fig:nn_dr} disaggregates performance to the class level, revealing
trade-offs the macro-F1 aggregate conceals.

\textbf{PCA} achieves the highest macro-F1 of \textbf{0.6951}, $+0.0148$
over the original baseline, supporting \textbf{H2}.
The class-level heatmap (\figref{fig:nn_dr}, right) shows where the gain
comes from: CT4 ($+0.040$), CT6 ($+0.044$), and CT3 ($+0.023$) ---
precisely the minority classes most burdened by correlated soil-indicator
noise that PCA filters out.
Balanced accuracy also improves ($0.664$ vs.\ $0.651$), consistent with
better minority-class recovery.

\textbf{ICA} (macro-F1 0.6785, $-0.002$) is broadly neutral but inflicts
localised losses on CT1 ($-0.013$) and CT6 ($-0.017$): independent
components preserve non-Gaussianity but shed the correlated variance that
benefits the two most common cover types.

\textbf{RP} (macro-F1 0.6670, $-0.013$) shows the most concentrated damage:
CT6 ($-0.073$) and CT7 ($-0.028$) collapse --- both classes occupy tight
local-geometry boundaries that random projections distort most severely.
The heatmap makes explicit that RP's macro-level deficit is not a uniform
penalty: it is a localised failure on the ecologically rarest classes.

Training time is modestly faster for reduced representations
($\sim 30\,\%$ reduction for PCA and RP), confirming DR can improve
both performance and compute simultaneously.

\begin{table}[!t]
\caption{NN after DR --- test performance and training time.
  \textbf{Best} bolded. DR fit on train split only (leakage-safe).}
\label{tab:nn_dr}\backto{sec:p4}
\centering\small
\rowcolors{2}{denimrow}{white}
\begin{tabular}{lrrrrr}
\toprule
\rowcolor{denimdark}
\textcolor{white}{\textbf{Features}} & \textcolor{white}{\textbf{dim}} & \textcolor{white}{\textbf{Macro-F1}} &
\textcolor{white}{\textbf{Accuracy}} & \textcolor{white}{\textbf{Bal.~Acc.}} & \textcolor{white}{\textbf{Train~(s)}} \\
\midrule
Original    & 54  & 0.6803 & 0.7940 & 0.6508 & 5.6 \\
PCA-30      & 30  & \textbf{0.6951} & \textbf{0.8000} & \textbf{0.6642} & \textbf{7.0} \\
ICA-20      & 20  & 0.6785 & 0.7875 & 0.6359 & 7.2 \\
RP-30       & 30  & 0.6670 & 0.7883 & 0.6339 & 6.9 \\
\bottomrule
\end{tabular}
\end{table}

\begin{figure}[!t]
  \centering
  \includegraphics[width=\linewidth]{./results/figures/fig_nn_dr_comparison.pdf}
  \caption{\emph{Left}: per-class F1 scores across all four feature spaces.
  PCA consistently helps the under-represented classes: CT4 (Cottonwood, $+$0.040),
  CT6 (Douglas-fir, $+$0.044), and CT3 (Ponderosa, $+$0.023) benefit most.
  \emph{Right}: $\Delta$F1 vs.\ original (red $=$ gain, blue $=$ loss).
  RP's largest damage lands on CT6 ($-$0.073) and CT7 ($-$0.028) ---
  both minority classes that rely on tight local-geometry boundaries
  that random projections distort.
  ICA is broadly neutral with slight losses on CT1 and CT6.
  The heatmap reveals that macro-F1 aggregates hide important class-level
  trade-offs not visible in \tabref{tab:nn_dr}.}
  \label{fig:nn_dr}\backto{sec:p4}
\end{figure}

% ── Section VII: Extra Credit — UMAP ─────────────────────────────────────────
\section{Extra Credit: UMAP}\label{sec:umap}

UMAP~\cite{mcinnes2018umap} ($n_{\text{neighbors}}{=}30$,
$\text{min\_dist}{=}0.1$, seed~7641) provides a 2D visualization for
comparison with PCA-2 (\figref{fig:2d}).
PCA-2 places CT1 and CT2 in overlapping elongated clouds
(the two dominant PCs capture elevation/hillshade variance that separates
most points but does not isolate ecological labels).
UMAP reveals sharper cluster topology: CT1 and CT2 form two distinct arcs
with visible gap, and minority classes CT4, CT5, CT7 cluster into tight
islands invisible in PCA-2 space.
This non-linear manifold structure confirms that the Covertype feature space
has local neighbourhood geometry beyond what linear projections capture.
Note: UMAP does not support standard transductive train/test transform; it
was used for visualization only and not for downstream NN evaluation.

\begin{figure}[!t]
  \centering
  \includegraphics[width=\linewidth]{./results/figures/fig_2d_visualization.pdf}
  \caption{2D projections of the 20k Covertype sample (3{,}000-point subsample).
  \emph{Top row}: PCA-2, coloured by Cover~Type (\emph{left}) and K-Means cluster
  (\emph{right}). CT1/CT2 form overlapping clouds; PCA captures elevation variance
  but not ecological labels.
  \emph{Bottom row}: UMAP-2 (extra credit; $n_{\text{neighbors}}{=}30$,
  $\text{min\_dist}{=}0.1$, seed~7641). Minority classes CT4, CT5, CT7
  form tight isolated islands in UMAP space that are invisible in PCA-2.
  K-Means cluster assignments (\emph{bottom right}) cut across these UMAP islands,
  showing elevation-driven clusters do not align with the non-linear
  manifold topology UMAP exposes.}
  \label{fig:2d}\backto{sec:umap}
\end{figure}

% ── Section VIII: Discussion ──────────────────────────────────────────────────
\section{Discussion}\label{sec:disc}

\subsection{Hypothesis review}\label{sec:disc:hyp}
\textbf{H1 --- supported, but the story is more interesting than the verdict.}
The challenge: does Covertype have natural cluster structure?
It does --- but not along ecological lines.
Both K-Means and GMM with $k{=}7$ yield ARI~$<0.06$, far below any
meaningful label-alignment threshold.
The resolution is that elevation and wilderness-area boundaries
\emph{do} produce real, stable clusters --- they simply do not track
ecological labels.
Clusters capture terrain, not taxonomy.
This is not a failure of the algorithms; it is a finding about the data.

\textbf{H2 --- supported, with an unexpected upside.}
The concern was whether 54-to-30 compression would hurt supervised learning.
The finding: PCA compression \emph{helped} (macro-F1 $+0.015$,
\tabref{tab:nn_dr} in \secref{sec:p4}).
PCA removes the correlated Hillshade columns and redundant soil indicators
that add noise without adding class-relevant signal, acting as implicit
regularisation.
The gain did not transfer to ICA or RP, which confirms that variance-preserving
compression --- not just dimensionality reduction in general --- is the mechanism.

\textbf{H3 --- partially supported; the RP sub-claim was falsified.}
PCA modestly improves K-Means silhouette (0.131 $\to$ 0.133), as predicted,
but RP gives the largest silhouette improvement (0.148) by aggressively
de-correlating the 44 sparse binary indicators.
The seed-variability claim was wrong: seed-to-seed variance of the \emph{mean}
ratio is 0.0045 --- negligible.
The real RP challenge is individual-pair distortion: ${\approx}20\%$ of pairs
fall outside $\pm 20\%$ of their original distance (\figref{fig:rp_dist}),
which is enough to shift centroids and widen covariance estimates in ways
that matter for both clustering and downstream NN.

\subsection{Feature geometry and intrinsic dimensionality}\label{sec:disc:efrank}
The effective rank of $\approx 5.93$ mentioned in \secref{sec:p2:pca} is
not just a PCA diagnostic --- it is the load-bearing number for interpreting
why \partref{4}{sec:p4} results look the way they do.
The ten continuous cartographic features collapse to six independent axes:
a Solar-Azimuth axis (PC1, 25.8\%), a Terrain-Accessibility axis (PC2,
21.6\%), and a Moisture-Proximity axis (PC3, 17.3\%).
The remaining 35.2\% of variance is distributed across seven further PCs
that mostly encode binary indicator noise (wilderness area and soil type
are the 44-dimensional block that contributes low-variance dimensions
after one-hot expansion).
This structure explains why PCA at $n{=}30$ \emph{improves} downstream MLP
performance ($+0.015$ Macro-F1, \tabref{tab:nn_dr}).
Of the 54 total dimensions, the retained 30 PCs capture 99.5\% of variance;
the discarded 24 (PCs 31--54) account for the remaining 0.5\%, which is
dominated by the sparse binary indicator block after standardisation.
These 24 dimensions add noise without adding signal --- PCA's truncation
effectively de-noises the binary block while preserving all cartographic
information.
ICA's instability ($|\kappa|$ seed std $= 4.58$) follows from the same
structure: FastICA is maximising non-Gaussianity across whitened inputs, but
the binary block's marginals are already maximally non-Gaussian by
construction (Bernoulli), so ICA has too many equally valid rotations to
choose from and different seeds converge to different orderings.
RP's individual-pair distortion (${\approx}20\%$ of pairs outside $\pm 20\%$)
is consistent with an effective rank of 6: random projections from 54 to 30
dimensions satisfy the Johnson--Lindenstrauss lemma
\emph{in expectation} for a 6-dimensional manifold, but individual pairs
that happen to align with the discarded 48-dimensional complement can
be severely distorted.

\textbf{Method assumptions and their effects.}
K-Means assumes compact spherical clusters; Covertype's imbalanced,
elongated class manifolds violate this assumption, producing low silhouette
even at the ``optimal'' $k{=}7$.
GMM with full covariance is a better geometric fit for Covertype but
requires more data per component; with only $\approx 2,860$ points per
component at $n{=}7$, the 54-dimensional covariance estimate is noisy,
as reflected in the negative GMM silhouette in original space.
After DR to 30 dimensions, GMM+PCA achieves the highest ARI (0.117),
suggesting the covariance estimation improves substantially with reduced
dimensionality.

ICA's instability (seed $|\kappa|$ std~$=$~4.58) is a practical limitation:
different runs yield different component orderings, making reproducibility
contingent on a fixed seed.
RP's strong distance preservation (mean ratio $\approx 1.0$) is theoretically
expected but the 12\,\% per-pair spread is sufficient to disrupt K-Means
centroid geometry, explaining the near-zero silhouette improvement.

\subsection{R sanity check}\label{sec:disc:rsanity}
The \texttt{rsanity/sanity\_check\_UL.Rmd} cross-checks Python results using
native R implementations.
R \texttt{stats::kmeans} (k$=$7, nstart$=$10) on the same scaled 20k sample
gives silhouette~$=$~0.1322 vs.\ Python's 0.1306
($|\Delta|{=}0.0016{<}0.02$ --- confirmed: the same feature matrix and
preprocessing are in use).
R \texttt{prcomp} agrees with sklearn PCA to floating-point precision:
$n_{90\%}{=}10$, $n_{95\%}{=}14$, PC1-EVR~$=$~0.226 --- identical.
All 10 clustering JSON logs pass range-validity checks in
\texttt{rsanity/sanity\_check\_UL.pdf}.

\subsection{Confusion matrix structure and elevation-band hypothesis}
\label{sec:disc:cm}
The per-class F1 heatmap (\figref{fig:nn_dr}) carries a structural message
that closes the loop on the \partref{1}{sec:p1} clustering.
The dominant MLP error pairs replicate the same confusion structure across
all four feature spaces: CT1$\leftrightarrow$CT2 (0.226 of Spruce-Fir
misclassified as Lodgepole), CT3$\leftrightarrow$CT6 (Ponderosa--Douglas-fir
cross-errors 0.134 and 0.292), and CT5$\to$CT2 (52\% of Aspen as Lodgepole).
These three confusion clusters map onto the three elevation bands that
K-Means found in \partref{1}{sec:p1} --- \emph{high-elevation} (CT1/CT2/CT7),
\emph{mid-elevation dry} (CT3/CT6), and \emph{transition-zone minority}
(CT4/CT5) --- confirming that the bottleneck is not the choice of
learner but the information content at those class boundaries.
PCA's improvement on CT3 ($+0.061$ recall) and CT4 ($+0.053$) reflects
DR removing the correlated hillshade noise that blurs the Ponderosa--
Douglas-fir boundary; the persistence of CT1$\leftrightarrow$CT2 confusion
even after PCA confirms a genuinely label-ambiguous co-occurrence zone
that no DR method can resolve.

\textbf{Limitations and evidence-based next steps.}
The 20k subsample preserves proportions but the rarest class (CT4, 94 instances)
has too few points for stable covariance estimation or silhouette measurement;
the $+0.053$ PCA recall gain for CT4 should be interpreted as a directional
finding, not a precise estimate.
Scaling only continuous columns is a deliberate choice but gives binary
indicators lower variance-influence in PCA; an alternative mixed-type
distance (Gower distance) might improve K-Means alignment by giving the
soil-type binary block appropriate weight in the Euclidean geometry.
ICA convergence warnings at 20 components suggest non-Gaussianity is
approaching the noise floor; reducing to 15 components and comparing
downstream NN performance would test whether the last 5 components
are contributing noise or signal.
The most valuable follow-up experiment grounded in the evidence: apply
\textbf{cost-sensitive reweighting} or SMOTE oversampling on CT4 and CT5
before training the downstream NN, then measure whether the PCA-preprocessed
feature space retains its $+0.053$ recall advantage for CT4 under
balanced training.
If it does, PCA's structural contribution (noise removal from the binary block)
is separable from the class-imbalance problem, and the two fixes are additive.
If it does not, the apparent PCA gain may partly reflect the rare-class
instability identified above rather than genuine geometric improvement.

% ── Section IX: Conclusion ────────────────────────────────────────────────────
\section{Conclusion}\label{sec:conclusion}

Three questions framed this study; three answers close it.

\textbf{Does Covertype have natural cluster structure?}
Yes, but not along ecological lines.
Elevation and wilderness-area geometry produces ARI~$\leq 0.052$ against
Cover~Type labels --- the terrain is real and partitions cleanly,
but the taxonomy is richer than the terrain.
Unsupervised methods are not wrong here; they are answering a different
question than the label would suggest.

\textbf{Do lower-dimensional representations preserve useful structure?}
PCA does, and it does better than the original.
At $n{=}30$ (99.5\% variance), PCA yields macro-F1~$=$~0.6951
($+0.015$ vs.\ original), better balanced accuracy, and faster training ---
confirming H2.
ICA captures genuine non-Gaussianity (mean $|\kappa|{=}6.05$) but is
unstable across seeds and does not improve downstream performance.
RP is theoretically sound but ${\approx}20\%$ of individual pairs are
distorted by $>{20\%}$, enough to degrade minority-class boundaries
($-0.013$ Macro-F1).

\textbf{What should a practitioner take away?}
Apply PCA preprocessing before any distance-sensitive method on Covertype.
Do not treat unsupervised clustering as a substitute for labelled training
on this problem --- the terrain structure is informative but insufficient.
UMAP reveals non-linear structure (tight minority-class islands invisible
in PCA-2 space) that semi-supervised or label-aware DR could exploit in
future work.

\textbf{The deeper finding, worth naming explicitly.}
A 581{,}012-row ecological classification dataset
--- collected to model which \emph{biological species} grows where ---
turns out to have an intrinsic dimensionality of only $\approx6$,
and all six independent axes are \emph{cartographic}: elevation, solar
exposure direction, terrain steepness, road proximity, and moisture
proximity.
No biological variable is in the feature set; the ecology is entirely
inferred from physical geography.
That the best supervised model (SL: SVM, macro-F1 $0.697$) still misses
one in every two Aspen (CT5) and Cottonwood (CT4) instances --- despite
perfect cartographic information --- is not a failure of the algorithm.
It is a signal that these transition-zone species are defined by
ecological processes (competition, disturbance, succession) that manifest
at finer spatial scales than the $30\,\mathrm{m} \times 30\,\mathrm{m}$
pixel and longer time scales than a single survey.
The unsupervised analysis makes this structural limit legible in a way
that the supervised accuracy number alone does not:
the clusters reveal terrain, and terrain reveals its limits.

% ── AI Use Statement ──────────────────────────────────────────────────────────
\section*{AI Use Statement}
Background reading on unsupervised learning methods --- the geometric
intuition behind K-Means vs.\ EM/GMM covariance structure, the whitening
rationale for FastICA, and the Johnson--Lindenstrauss guarantees for random
projections --- was sourced in part via Gemini (Google) queries about specific
algorithmic properties and edge cases (e.g.\ why ICA kurtosis can be
seed-sensitive on mixed continuous/binary inputs, and when RP distance
preservation degrades with sparse binary features); the quantitative claims
are independently computed from the UCI data and the cited primary literature.
Proofreading and grammar passes used Grammarly, which also helped
naturalise phrasing that reads awkwardly for readers more accustomed to
North American academic English idioms;
code debugging conversations ran through DeepSeek;
light refactoring of plotting and table-emission utilities used Visual Studio
Code Copilot.
All hypotheses were locked in \texttt{notes/HYPOTHESIS\_UL.md} before
running any \partref{1}{sec:p1}--4 experiments; every reported number was re-produced by
the student's own runs of the underlying code on the student's own hardware;
the \LaTeX{} prose was authored by the student with the LLM assistance
described above and revised in full.
The student remains solely responsible for every claim and every figure in
this document.

% ── References ────────────────────────────────────────────────────────────────
\begin{thebibliography}{9}

\bibitem{jolliffe2002pca}
I.~T. Jolliffe,
\textit{Principal Component Analysis}, 2nd~ed.
New York: Springer, 2002.
ISBN 978-0-387-95442-4.

\bibitem{hyvarinen2000ica}
A.~Hyv\"{a}rinen and E.~Oja,
``Independent component analysis: algorithms and applications,''
\textit{Neural Networks}, vol.~13, no.~4--5, pp.~411--430, 2000.
\url{https://doi.org/10.1016/S0893-6080(00)00026-5}

\bibitem{johnson1984jl}
W.~B. Johnson and J.~Lindenstrauss,
``Extensions of Lipschitz mappings into a Hilbert space,''
in \textit{Proc. Conf. Modern Analysis and Probability},
pp.~189--206, 1984.
\url{https://doi.org/10.1090/conm/026/737400}

\bibitem{mcinnes2018umap}
L.~McInnes, J.~Healy, and J.~Melville,
``UMAP: Uniform Manifold Approximation and Projection for Dimension Reduction,''
\textit{arXiv:1802.03426}, 2018.
\url{https://arxiv.org/abs/1802.03426}

\bibitem{sklearn2024}
F.~Pedregosa \textit{et al.},
``Scikit-learn: Machine learning in Python,''
\textit{J. Mach. Learn. Res.}, vol.~12, pp.~2825--2830, 2011.
\url{https://jmlr.org/papers/v12/pedregosa11a.html}

\bibitem{rousseeuw1987sil}
P.~J. Rousseeuw,
``Silhouettes: a graphical aid to the interpretation and validation of cluster analysis,''
\textit{J. Comput. Appl. Math.}, vol.~20, pp.~53--65, 1987.
\url{https://doi.org/10.1016/0377-0427(87)90125-7}

\end{thebibliography}

\end{document}

